\documentclass{report}
\usepackage{luatexja} % LuaTeXで日本語を使うためのパッケージ
\usepackage{luatexja-fontspec} % LuaTeX用の日本語フォント設定

% --- 数学関連 ---
\usepackage{amsmath, amssymb, amsfonts, mathtools, bm, amsthm} % 基本的な数学パッケージ
\usepackage{type1cm, upgreek} % 数式フォントとギリシャ文字
\usepackage{physics, mhchem} % 物理や化学の記号や式の表記を簡単にする

% --- 表関連 ---
\usepackage{multirow, longtable, tabularx, array, colortbl, dcolumn, diagbox} % 表のレイアウトを柔軟にする
\usepackage{tablefootnote, truthtable} % 表中に注釈を追加、真理値表
\usepackage{tabularray} % 高度な表組みレイアウト

% --- グラフィック関連 ---
\usepackage{tikz, graphicx} % 図の描画と画像の挿入
\usepackage{background} % ウォーターマークの設定
\usepackage{caption, subcaption} % 図や表のキャプション設定
\usepackage{float, here} % 図や表の位置指定

% --- レイアウトとページ設定 ---
\usepackage{fancyhdr} % ページヘッダー、フッター、余白の設定
\usepackage[top = 20truemm, bottom = 20truemm, left = 20truemm, right = 20truemm]{geometry}
\usepackage{fancybox, ascmac} % ボックスのデザイン

% --- 色とスタイル ---
\usepackage{xcolor, color, colortbl, tcolorbox} % 色とカラーボックス
\usepackage{listings, jvlisting} % コードの色付けとフォーマット

% --- 参考文献関連 ---
\usepackage{biblatex, usebib} % 参考文献の管理と挿入
\usepackage{url, hyperref} % URLとリンクの設定

% --- その他の便利なパッケージ ---
\usepackage{footmisc} % 脚注のカスタマイズ
\usepackage{multicol} % 複数段組
\usepackage{comment} % コメントアウトの拡張
\usepackage{siunitx} % 単位の表記
\usepackage{docmute}
% \usepackage{appendix}
% --- tcolorboxとtikzの設定 ---
\tcbuselibrary{theorems, breakable} % 定理のボックスと改ページ設定
\usetikzlibrary{decorations.markings, arrows.meta, calc} % tikzの装飾や矢印の設定

% --- 定理スタイルと数式設定 ---
\theoremstyle{definition} % 定義スタイル
\numberwithin{equation}{section} % 式番号をサブセクション単位でリセット

% --- hyperrefの設定 ---
\hypersetup{
  setpagesize = false,
  bookmarks = true,
  bookmarksdepth = tocdepth,
  bookmarksnumbered = true,
  colorlinks = false,
  pdftitle = {}, % PDFタイトル
  pdfsubject = {}, % PDFサブジェクト
  pdfauthor = {}, % PDF作者
  pdfkeywords = {} % PDFキーワード
}

% --- siunitxの設定 ---
\sisetup{
  table-format = 1.5, % 小数点以下の桁数
  table-number-alignment = center, % 数値の中央揃え
}

% --- 透かし画像の設定 --- 
\backgroundsetup{
  scale=0.5,                       % 画像のスケール
  % color=black,                   % 画像の色（透かし用に半透明が推奨）
  opacity=0.2,                   % 透かしの透明度（0が完全透明、1が完全不透明）
  angle=0,                       % 画像の角度
  position = current page.south east,  % ページの右下
  hshift=-6cm, % 右方向へのシフト（負の値で内側に移動）
  vshift=5cm,  % 上方向へのシフト（正の値で内側に移動）
  contents={\includegraphics{./fig/appilogo-circular-full.png}} % 画像のパス
}

% --- その他の設定 ---
\allowdisplaybreaks % 数式の途中改ページ許可
\newcolumntype{t}{!{\vrule width 0.1pt}} % 新しいカラムタイプ
\newcolumntype{b}{!{\vrule width 1.5pt}} % 太いカラム
\UseTblrLibrary{amsmath, booktabs, counter, diagbox, functional, hook, html, nameref, siunitx, varwidth, zref} % tabularrayのライブラリ
\setlength{\columnseprule}{0.4pt} % カラム区切り線の太さ
\captionsetup[figure]{font = bf} % 図のキャプションの太字設定
\captionsetup[table]{font = bf} % 表のキャプションの太字設定
\captionsetup[lstlisting]{font = bf} % コードのキャプションの太字設定
\captionsetup[subfigure]{font = bf, labelformat = simple} % サブ図のキャプション設定
\setcounter{secnumdepth}{4} % セクションの深さ設定
\newcolumntype{d}{D{.}{.}{5}} % 数値のカラム
\newcolumntype{M}[1]{>{\centering\arraybackslash}m{#1}} % センター揃えのカラム
\DeclareMathOperator{\diag}{diag}
\everymath{\displaystyle} % 数式のスタイル
\newcommand{\inner}[2]{\left\langle #1, #2 \right\rangle}
\renewcommand{\figurename}{図}
\renewcommand{\i}{\mathrm{i}} % 複素数単位i
\renewcommand{\laplacian}{\grad^2} % ラプラシアンの記号
\renewcommand{\thesubfigure}{(\alph{subfigure})} % サブ図の番号形式
\newcommand{\m}[3]{\multicolumn{#1}{#2}{#3}} % マルチカラムのショートカット
\renewcommand{\r}[1]{\mathrm{#1}} % mathrmのショートカット
\newcommand{\e}{\mathrm{e}} % 自然対数の底e
\newcommand{\Ef}{E_{\mathrm{F}}} % フェルミエネルギー
\renewcommand{\c}{\si{\degreeCelsius}} % 摂氏記号
\renewcommand{\d}{\r{d}} % d記号
\renewcommand{\t}[1]{\texttt{#1}} % タイプライタフォント
\newcommand{\kb}{k_{\mathrm{B}}} % ボルツマン定数
% \renewcommand{\phi}{\varphi} % ϕをφに変更
\renewcommand{\epsilon}{\varepsilon}
\newcommand{\fullref}[1]{\textbf{\ref{#1} \nameref{#1}}}
\newcommand{\reff}[1]{\textbf{図\ref{#1}}} % 図参照のショートカット
\newcommand{\reft}[1]{\textbf{表\ref{#1}}} % 表参照のショートカット
\newcommand{\refe}[1]{\textbf{式\eqref{#1}}} % 式参照のショートカット
\newcommand{\refp}[1]{\textbf{コード\ref{#1}}} % コード参照のショートカット
\renewcommand{\lstlistingname}{コード} % コードリストの名前
\renewcommand{\theequation}{\thesection.\arabic{equation}} % 式番号の形式
\renewcommand{\footrulewidth}{0.4pt} % フッターの線
\newcommand{\mar}[1]{\textcircled{\scriptsize #1}} % 丸囲み文字
\newcommand{\combination}[2]{{}_{#1} \mathrm{C}_{#2}} % 組み合わせ
\newcommand{\thline}{\noalign{\hrule height 0.1pt}} % 細い横線
\newcommand{\bhline}{\noalign{\hrule height 1.5pt}} % 太い横線

% --- カスタム色定義 ---
\definecolor{burgundy}{rgb}{0.5, 0.0, 0.13} % バーガンディ色
\definecolor{charcoal}{rgb}{0.21, 0.27, 0.31} % チャコール色
\definecolor{forest}{rgb}{0.0, 0.35, 0} % 森の緑色

% --- カスタム定理環境の定義 ---
\newtcbtheorem[number within = chapter]{myexc}{練習問題}{
  fonttitle = \gtfamily\sffamily\bfseries\upshape,
  colframe = forest,
  colback = forest!2!white,
  rightrule = 1pt,
  leftrule = 1pt,
  bottomrule = 2pt,
  colbacktitle = forest,
  theorem style = standard,
  breakable,
  arc = 0pt,
}{exc-ref}
\newtcbtheorem[number within = chapter]{myprop}{命題}{
  fonttitle = \gtfamily\sffamily\bfseries\upshape,
  colframe = blue!50!black,
  colback = blue!50!black!2!white,
  rightrule = 1pt,
  leftrule = 1pt,
  bottomrule = 2pt,
  colbacktitle = blue!50!black,
  theorem style = standard,
  breakable,
  arc = 0pt
}{proposition-ref}
\newtcbtheorem[number within = chapter]{myrem}{注意}{
  fonttitle = \gtfamily\sffamily\bfseries\upshape,
  colframe = yellow!20!black,
  colback = yellow!50,
  rightrule = 1pt,
  leftrule = 1pt,
  bottomrule = 2pt,
  colbacktitle = yellow!20!black,
  theorem style = standard,
  breakable,
  arc = 0pt
}{remark-ref}
\newtcbtheorem[number within = chapter]{myex}{例題}{
  fonttitle = \gtfamily\sffamily\bfseries\upshape,
  colframe = black,
  colback = white,
  rightrule = 1pt,
  leftrule = 1pt,
  bottomrule = 2pt,
  colbacktitle = black,
  theorem style = standard,
  breakable,
  arc = 0pt
}{example-ref}
\newtcbtheorem[number within = chapter]{exc}{Requirement}{myexc}{exc-ref}
\newcommand{\rqref}[1]{{\bfseries\sffamily 練習問題 \ref{exc-ref:#1}}}
\newtcbtheorem[number within = chapter]{definition}{Definition}{mydef}{definition-ref}
\newcommand{\dfref}[1]{{\bfseries\sffamily 定義 \ref{definition-ref:#1}}}
\newtcbtheorem[number within = chapter]{prop}{命題}{myprop}{proposition-ref}
\newcommand{\prref}[1]{{\bfseries\sffamily 命題 \ref{proposition-ref:#1}}}
\newtcbtheorem[number within = chapter]{rem}{注意}{myrem}{remark-ref}
\newcommand{\rmref}[1]{{\bfseries\sffamily 注意 \ref{remark-ref:#1}}}
\newtcbtheorem[number within = chapter]{ex}{例題}{myex}{example-ref}
\newcommand{\exref}[1]{{\bfseries\sffamily 例題 \ref{example-ref:#1}}}
% --- 再定義コマンド ---
% \mathtoolsset{showonlyrefs=true} % 必要な式番号のみ表示
\pagestyle{fancy} % ヘッダー・フッターのスタイル設定
\chead{応用量子物性講義ノート} % 中央ヘッダー
% \rhead{}
\fancyhead[R]{\rightmark}
\renewcommand{\sectionmark}[1]{\markright{\thesection\ #1}}
\cfoot{\thepage} % 中央フッターにページ番号
\lhead{}
\rfoot{Yuto Masuda and Haruki Aoki} % 右フッターに名前
\setcounter{tocdepth}{4} % 目次の深さ
\makeatletter
\@addtoreset{equation}{section} % サブセクションごとに式番号をリセット
\makeatother

% --- メタ情報 ---
\title{応用量子物性講義ノート}
\date{更新日\today}
\author{Yuto Masuda and Haruki Aoki}

\begin{document}
角運動量演算子の行列表現を求める．使うのは交換関係
\begin{align}
  \qty[\hat{L}_i,\hat{L}_j] = \i\hbar\epsilon_{ijk}\hat{L}_k \label{commutation-of-am}
\end{align}
のみである．まず，
\begin{align}
  \qty[\hat{L}^2, \hat{L}_i] = 0
\end{align}
は簡単に確かめられる．よって，$\hat{L}^2$と$\hat{L}_z$は同時固有状態を持つ．同時固有状態を$\ket{\alpha\beta}$とし固有値を
\begin{align}
  \hat{L}^2\ket{\alpha,\beta} &= \alpha\ket{\alpha,\beta}\\
  \hat{L}_z\ket{\alpha,\beta} &= \beta\ket{\alpha,\beta}
\end{align}
と定める．

次に，\textbf{昇降演算子}(ladder operator)を
\begin{align}
  \hat{L}_{\pm} \equiv \hat{L}_x \pm \i \hat{L}_y
\end{align}
と定義する．
\begin{align}
  \hat{L}_z(\hat{L}_{+}\ket{\alpha,\beta}) = (\beta + \hbar)\hat{L}_{+}\ket{\alpha,\beta}
\end{align}
が計算により確かめられるため，$\hat{L}_{+}\ket{\alpha\beta}$は$\hat{L}_z$の固有値$\beta+\hbar$の固有状態である．一方，$[\hat{L}^2, \hat{L}_{\pm}]=0$であるため
昇降演算子は$L$の大きさを変化させない．変化するのは$\beta$のみである．また，
\begin{align}
  &\mel**{\psi}{\hat{L}^2 - \hat{L}_z^2 }{\psi}\\
  &= \mel**{\psi}{\hat{L}_x^2}{\psi} + \mel**{\psi}{\hat{L}_y^2}{\psi} \geq 0\\
  &\Rightarrow \ev{ \hat{L}_z^2 } \leq \ev{ \hat{L}^2 }
\end{align}
であるため，$\beta$には最大値$\beta_{\r{max}}$が存在する．よって，あるところで
\begin{align}
  \hat{L}_{+}\ket{\alpha,\beta_{\r{max}}} &= 0\\
  \hat{L}_{-}\ket{\alpha,\beta_{\r{min}}} &= 0
\end{align}
が成り立つ．
\begin{align}
  \hat{L}_{-}\hat{L}_{+} = \hat{L}^2 - \hat{L}_z^2 - \hbar \hat{L}_z
\end{align}
だから，
\begin{align}
  0 = (\alpha - \beta_{\r{max}}^2 - \hbar \beta_{\r{\max}})\ket{\alpha,\beta_{\r{max}}}
\end{align}
である．よって，
\begin{align}
  \alpha = \beta_{\r{max}}(\beta_{\r{max}} + \hbar)
\end{align}
が得られる．同様に
\begin{align}
  \alpha = \beta_{\r{min}}(\beta_{\r{min}} - \hbar)
\end{align}
である．この2式を比べると
\begin{align}
  \beta_{\r{max}} = \beta_{\r{min}}
\end{align}
がわかる．また$\hat{L}_z$の固有値は昇降演算子により$\hbar$の整数倍で増減するから$k$を自然数として
\begin{align}
  \beta_{\r{max}} = \beta_{\r{min}} + k\hbar
\end{align}
と書ける．よって，
\begin{align}
  \beta_{\r{max}} = \frac{k}{2}\hbar
\end{align}
である．つまり，$\hat{L}_z$の最大値は$\frac{\hbar}{2}$の自然数倍の値しかとらない．さらに，
\begin{align}
  \alpha = \hbar^2\qty(\frac{k}{2})\qty(\frac{k}{2} + 1)
\end{align}
である．ここで
\begin{align}
  l \equiv \frac{k}{2}
\end{align}
とすれば
\begin{align}
  \hat{L}^2\ket{l,m} &= \hbar^2 l(l+1)\ket{l,m} \label{abs-of-l}\\
  \hat{L}_z\ket{l,m} &= m\hbar\ket{l,m} \label{magnitude-of-lz}
\end{align}
と書ける．ここで，
\begin{align}
  l &= 0,\frac{1}{2},1,\frac{3}{2},\cdots\\
  m &= -l,-l+1,-l+2,\cdots,l-2,l-1,l
\end{align}
である．

次に，
\begin{align}
  \hat{L}_{+}\ket{l,m} = C_{+}\ket{l,m+1}
\end{align}
とおく．$\hat{L}_{+}^{\dagger} = \hat{L}_{-}$だから，
\begin{align}
  &\abs{C_{+}}^2 = \bra{l,m}\hat{L}_{-}\hat{L}_{+}\ket{l,m} = \qty[l(l+1) -m(m+1)]\hbar^2\\
  &\Rightarrow C_{+} = \hbar\sqrt{(l-m)(l+m+1)}
\end{align}
である．$\hat{L}_{-}$に関しても同様に計算すると，
\begin{align}
  \hat{L}_{+}\ket{l,m} = \hbar\sqrt{(l-m)(l+m+1)}\ket{l,m+1}\\
  \hat{L}_{-}\ket{l,m} = \hbar\sqrt{(l+m)(l-m+1)}\ket{l,m+1}
\end{align}
を得る．よって，上2式と\refe{abs-of-l}と\refe{magnitude-of-lz}を使えば$\hat{L}_z$と$\hat{L}^2$の行列表現が得られる．
\begin{align}
  &\hat{L}_z = \\
  &\begin{pmatrix} 
    \mel{0,0}{\hat{L}_z}{0,0} & \mel{0,0}{\hat{L}_z}{1/2,1/2} & \mel{0,0}{\hat{L}_z}{1/2,-1/2} & \mel{0,0}{\hat{L}_z}{1,1} & \mel{0,0}{\hat{L}_z}{1,0} & \cdots \\
    \mel{1/2,1/2}{\hat{L}_z}{0,0} & \mel{1/2,1/2}{\hat{L}_z}{1/2,1/2} & \mel{1/2,1/2}{\hat{L}_z}{1/2,-1/2} & \mel{1/2,1/2}{\hat{L}_z}{1,1} & \mel{1/2,1/2}{\hat{L}_z}{1,0} & \cdots \\
    \mel{1/2,-1/2}{\hat{L}_z}{0,0} & \mel{1/2,-1/2}{\hat{L}_z}{1/2,1/2} & \mel{1/2,-1/2}{\hat{L}_z}{1/2,-1/2} & \mel{1/2,-1/2}{\hat{L}_z}{1,1} & \mel{1/2,-1/2}{\hat{L}_z}{1,0} & \cdots \\
    \mel{1,1}{\hat{L}_z}{0,0} & \mel{1,1}{\hat{L}_z}{1/2,1/2} & \mel{1,1}{\hat{L}_z}{1/2,-1/2} & \mel{1,1}{\hat{L}_z}{1,1} & \mel{1,1}{\hat{L}_z}{1,0} & \cdots \\
    \mel{1,0}{\hat{L}_z}{0,0} & \mel{1,0}{\hat{L}_z}{1/2,1/2} & \mel{1,0}{\hat{L}_z}{1/2,-1/2} & \mel{1,0}{\hat{L}_z}{1,1} & \mel{1,0}{\hat{L}_z}{1,0} & \cdots \\
    \vdots & \vdots & \vdots & \vdots  
  \end{pmatrix} \\
  &=
  \begin{pmatrix}
    0 & 0 & 0 & 0 & 0 & \cdots\\
    0 & \frac{\hbar}{2} &  &  &  & \\
    0  & & -\frac{\hbar}{2} &  & & \\
    \vdots & & & \hbar & & \cdots \\
    \vdots & & & & 0 &  \cdots\\
    \vdots & & & & & -\hbar 
  \end{pmatrix}
\end{align}
\begin{align}
  \hat{L}^2 =
  \begin{pmatrix}
    0 & 0 & & & & \\
    0 & \frac{3}{4}\hbar^2 & & & & \\
    & & \frac{3}{4}\hbar^2 & & & \\
    & & & 2\hbar^2 & & \\
    & & & & 2\hbar^2 & \\
    & & & & & 2\hbar^2 \\
  \end{pmatrix}
\end{align}
注目すべきは$l$ごとにブロック状になっていることである．$l=1/2$の部分だけ取り出せば
\begin{align}
  \hat{L}_z = \frac{\hbar}{2}
  \begin{pmatrix}
    1 & 0\\
    0 & -1
  \end{pmatrix}
\end{align}
が得られる．
\end{document}